% JHEP 体裁の骨格。jheppub.sty を同じディレクトリに置く（references/latex-conventions.md）。
% jheppub を使わない場合は \usepackage{jheppub} を外し、\abstract{...} を abstract 環境に、
% \affiliation・\emailAdd・\keywords・\preprint を通常の article の書き方に直す。
\documentclass[a4paper,11pt]{article}
\usepackage{jheppub}
\usepackage{amsmath,amssymb,amsthm}
\usepackage{booktabs}

% 定理環境は、証明のある主張と主要な予想にだけ使う。
\newtheorem{proposition}{Proposition}[section]
\newtheorem{conjecture}[proposition]{Conjecture}

\title{Title of the paper}
\author{First Author}
\affiliation{Institute, Address}
\emailAdd{author@example.org}

% 1 段落。引用・ディスプレイ式なし。定式化したこと、予想、確かめた範囲を書き分ける。
\abstract{We formulate ... We conjecture that ... We show that, for ..., ...}

\begin{document}
\maketitle

\section{Introduction}
% 文脈 -> 未解決点 -> この論文の寄与（散文）-> 既知の結果との境界 -> 構成の段落。

The paper is organized as follows. In section~\ref{sec:review} we review ...

\section{Review and conventions}\label{sec:review}
% 標準的な定義は引用で済ませる。固有の記法と規約はここで 1 回だけ定義する。

\section{Main results}\label{sec:main}

\section{Examples}\label{sec:examples}
% 計算による確認は、範囲を明記した 1〜2 文と付録の表への参照にとどめる。

\section{Discussion}\label{sec:discussion}

\acknowledgments
% 議論の相手、グラント。AI の使用の開示（文面は著者が確定する）。
% The author used Claude (Anthropic) for computations, literature search and drafting.
% All results, references and text were checked by the author, who takes full responsibility
% for the content.

\appendix
\section{Conventions}\label{app:conventions}
\section{Tables}\label{app:data}

\bibliographystyle{JHEP}
\bibliography{refs}
\end{document}
